\answer{ex:05:1}{$\ell\approx17\um$; $\approx100$~jours: la diffusion générale dans le cerveau est assurée par la ramification du fil, la diffusion moléculaire ne fait qu'étaler chaque site de libération sur $\sim20\um$.}
\answer{ex:05:2}{$r=ab/(1+G)$, $S=1/(1+G)$ exactement; $+1.7\,\%$ au lieu de $+20\,\%$.}
\answer{ex:05:3}{$T^*=1.21\,\text{s}$ pour $G=2$ et $0.19\,\text{s}$ pour $G=9$: en pratique $G\sim2$--$4$, soit une réduction d'un facteur $3$ à $5$.}
\answer{ex:05:4}{$\mathrm{CV}=1/\sqrt{2\lambda\tau_c}\approx9\cdot10^{-4}$ pour une fréquence totale $\lambda$; un seul neurone demande $\tau_c=12.5\,\text{s}$.}
\answer{ex:05:5}{$c\approx0.40$ dans les deux cas; $\mathrm{CV}=0.050$ contre $0.158$, soit un facteur $\sqrt{10}$. Les fréquences des neurones restent proportionnelles aux $a_i$: seule la somme est régulée.}
\answer{ex:05:6}{$N_1=\overline{A\vee B}$, $N_2=\overline{A\vee N_1}$, $N_3=\overline{B\vee N_1}$, $N_4=\overline{N_2\vee N_3}$, XOR $=N_5=\overline{N_4}$; chaque commutation prend $\tau_c\ln2$, le chemin de $4$ portes $2.8\,\text{s}$ par XOR.}
\answer{ex:05:7}{(a)~$\tau_\gamma=6/\ln3\approx5.5\,\text{s}$. (b)~$\bar D<4\,\text{s}$ contre $D<2.2\,\text{s}$: l'imprévisibilité du délai rend plus patient.}
\answer{ex:05:8}{$k_2=1/\tau_d=10\,\text{s}^{-1}$, $k_1=1/\tau_d-1/\tau_\gamma=9.5\,\text{s}^{-1}$; $\tau_\gamma=1/(k_2-mk_1)$: doublement pour $+2.6\,\%$ (et pour $+5.3\,\%$ l'horizon devient infini).}
